\chapter{第七单元}

\Needspace{12\baselineskip}

\section{第七单元素养评价（A）}

\Needspace{8\baselineskip}

\pdfbookmark[2]{原书第57页}{orig-57}

\subsection*{原书第57页}

% p057_s01
\BookQuestion{一、1}
计算小数加减法时，要先把（\hspace{1.8em}）对齐，也就是把（\hspace{1.8em}）对齐。\par
\begin{BookAnswer}
小数加减法先把“小数点”对齐，也就是把“相同数位”对齐；这样才能把相同计数单位的个数相加减。计算要从低位向高位，满十向前一位进1，不够减向前一位退1；退1在后一位就表示10个这样的计数单位。
\end{BookAnswer}

% p057_s02
\BookQuestion{一、2}
填空：\(6.3+4.9+13.7=4.9+(\square+\square)\)；\allowbreak{}\(24.8-5.6-4.4=24.8-(\square\mathbin{\bigcirc}\square)\)。\par
\begin{BookAnswer}
第一式同时用加法交换律与结合律：\(6.3+4.9+13.7=4.9+(6.3+13.7)\)，填6.3、13.7。第二式用减法性质：\(24.8-5.6-4.4=24.8-(5.6+4.4)\)，填5.6、“+”、4.4。
\end{BookAnswer}

% p057_s03
\BookQuestion{一、3}
两个加数的和是\(36.8\)，一个加数增加\(5.9\)，现在的和是（\hspace{1.8em}）。\par
\begin{BookAnswer}
另一个加数不变，一个加数增加5.9，和也增加5.9。把5.9凑成6，多加的0.1再减回：\[\begin{aligned}&36.8+5.9\\={}&36.8+6-0.1\\={}&42.8-0.1\\={}&42.7\end{aligned}\]填42.7。
\end{BookAnswer}

% p057_s04
\BookQuestion{一、4}
20个\(0.1\)和20个\(0.01\)，它们的和是（\hspace{1.8em}），差是（\hspace{1.8em}）。\par
\begin{BookAnswer}
20个0.1是2，20个0.01是0.2；和为\(2+0.2=2.2\)，差为\(2-0.2=1.8\)。
\end{BookAnswer}

% p057_s05
\BookQuestion{一、5（1）}
按规律填数：（1）\(3.2,4.0,4.8,(\quad),(\quad)\)。\par
\begin{BookAnswer}
每次增加0.8：\(3.2,4.0,4.8,5.6,6.4\)。所以填5.6、6.4。
\end{BookAnswer}

% p057_s06
\BookQuestion{一、5（2）}
按规律填数：（2）\(10.1,8.9,7.7,(\quad),(\quad)\)。\par
\begin{BookAnswer}
每次减少1.2：\(10.1,8.9,7.7,6.5,5.3\)。填6.5、5.3。
\end{BookAnswer}

% p057_s07
\BookQuestion{一、5（3）}
按规律填数：（3）\(3.5,9.5,15.5,(\quad),27.5\)。\par
\begin{BookAnswer}
每次增加6：\(3.5,9.5,15.5,21.5,27.5\)。填21.5。
\end{BookAnswer}

% p057_s08
\BookQuestion{一、6}
\(0.46\)扩大到原来的10倍，比原数大（\hspace{1.8em}）。\par
\begin{BookAnswer}
扩大到原来的10倍，得到4.6；比原数大的部分为\(4.6-0.46=4.14\)。也可理解为比原来增加9个0.46。
\end{BookAnswer}

% p057_s09
\BookQuestion{一、7}
甲、乙两数的和是\(26.3\)，甲数是\(8.6\)，乙数是（\hspace{1.8em}）。\par
\begin{BookAnswer}
一个加数=和−另一个加数，所以乙数为\(26.3-8.6=17.7\)。
\end{BookAnswer}

% p057_s10
\BookQuestion{一、8}
甲数是\(9.2\)，比乙数多\(3.7\)，乙数是（\hspace{1.8em}）。\par
\begin{BookAnswer}
甲比乙多3.7，求较小的乙数用减法：\(9.2-3.7=5.5\)。
\end{BookAnswer}

% p057_s11
\BookQuestion{一、9}
\(4.8\)的计数单位是（\hspace{1.8em}），它有（\hspace{1.8em}）个这样的计数单位，再添（\hspace{1.8em}）个这样的计数单位就是5。\par
\begin{BookAnswer}
4.8是一位小数，计数单位是0.1。4中有40个0.1，0.8中有8个0.1，合起来有48个0.1。\(5-4.8=0.2\)，还要添2个0.1。填0.1、48、2。
\end{BookAnswer}

% p057_s12
\BookQuestion{一、10（1）}
在圆圈里填上“＞”“＜”或“＝”：\(5.6+4.4\ \bigcirc\ 12-1.7\)。\par
\begin{BookAnswer}
左边\(5.6+4.4=10\)，右边\(12-1.7=10.3\)；\allowbreak{}\(10<10.3\)，填“<”。
\end{BookAnswer}

% p057_s13
\BookQuestion{一、10（2）}
在圆圈里填上“＞”“＜”或“＝”：\(1.08+5.99\ \bigcirc\ 7.08\)。\par
\begin{BookAnswer}
左边\(1.08+5.99=7.07\)，\allowbreak{}\(7.07<7.08\)，填“<”。
\end{BookAnswer}

% p057_s14
\BookQuestion{一、10（3）}
在圆圈里填上“＞”“＜”或“＝”：\(13.5-6\ \bigcirc\ 4.8+1.9\)。\par
\begin{BookAnswer}
左边\(13.5-6=7.5\)，右边\(4.8+1.9=6.7\)，填“>”。
\end{BookAnswer}

% p057_s15
\BookQuestion{一、10（4）}
在圆圈里填上“＞”“＜”或“＝”：\(7.3+5.9\ \bigcirc\ 8.2+4.7\)。\par
\begin{BookAnswer}
左边\(7.3+5.9=13.2\)，右边\(8.2+4.7=12.9\)，填“>”。
\end{BookAnswer}

% p057_s16
\BookQuestion{一、11}
一根木条长\(18.2\)厘米，两段拼接重叠\(3.6\)厘米，拼好总长29厘米，另一根木条长（\hspace{1.8em}）厘米。\par
\begin{BookAnswer}
两根木条原长的和等于拼好后的总长加重叠部分。另一根长度：\[\begin{aligned}&29+3.6-18.2\\={}&32.6-18.2\\={}&14.4\text{(厘米)}\end{aligned}\]填14.4。重叠部分在两根原长中算两次，在拼好总长中算一次。
\end{BookAnswer}

\Needspace{8\baselineskip}

\pdfbookmark[2]{原书第58页}{orig-58}

\subsection*{原书第58页}

% p058_s01
\BookQuestion{一、12}
今日气温\(19.5\,{}^\circ\mathrm C\)，明日比今日高\(3.8\,{}^\circ\mathrm C\)，明日气温（\hspace{1.8em}）\({}^\circ\mathrm C\)。\par
\begin{BookAnswer}
明日比今日高3.8℃，所以\(19.5+3.8=23.3\)℃，填23.3。
\end{BookAnswer}

% p058_s02
\BookQuestion{二、1}
一个数加\(9.5\)等于30，这个数是（\hspace{1.8em}）。A.\(39.5\)；B.\(21.5\)；C.\(20.5\)。\par
\begin{BookAnswer}
一个加数等于和减另一个加数。把9.5凑成10，多减的0.5再加回：\[\begin{aligned}&30-9.5\\={}&30-10+0.5\\={}&20+0.5\\={}&20.5\end{aligned}\]所以选C。
\end{BookAnswer}

% p058_s03
\BookQuestion{二、2}
5元3角＋8角7分＝（\hspace{1.8em}）。A.\(6.17\)元；B.14元；C.\(6.1\)元。\par
\begin{BookAnswer}
先把金额都化成元：\[\begin{aligned}5\text{元}3\text{角}&=5.30\text{元}\\8\text{角}7\text{分}&=0.87\text{元}\end{aligned}\]再算\(5.30+0.87=6.17\)元，选A。
\end{BookAnswer}

% p058_s04
\BookQuestion{二、3}
一个一位小数加上一个两位小数，和是（\hspace{1.8em}）位小数。A.一；B.二；C.三。\par
\begin{BookAnswer}
先把一位小数的末尾添0，改写成两位小数，再把相同数位相加。若按竖式计算后暂不化简、保留两位的写法，可选B；题目未说明该约定，严格说结论不唯一。\par 须区分数值与写法：小数末尾添0或去0不改变大小，如\(1.2+0.30=1.50=1.5\)。若把结果去掉小数末尾的0，还可能得到整数，例如\(1.2+1.80=3.00=3\)。因此按最简写法，和可能是整数、一位小数或两位小数；而按竖式暂不化简的两位写法，才统一填B。
\end{BookAnswer}

% p058_s05
\BookQuestion{二、4}
被减数不变，减数增加\(3.6\)，差（\hspace{1.8em}）。A.增加\(3.6\)；B.减少\(3.6\)；C.不变。\par
\begin{BookAnswer}
被减数不变，减数增加3.6，差相应减少3.6，选B。
\end{BookAnswer}

% p058_s06
\BookQuestion{二、5}
\(76.54+9.99=(\quad)\)。A.\(76.54+10+0.01\)；B.\(76.54+10-0.01\)；C.\(76.54-10+0.01\)。\par
\begin{BookAnswer}
\(9.99=10-0.01\)，因此\(76.54+9.99=76.54+10-0.01\)，选B。
\end{BookAnswer}

% p058_s07
\BookQuestion{二、6}
四年级百米跑步比赛，小明用时\(17.8\)秒，比小雨慢\(0.9\)秒，小雨百米用时（\hspace{1.8em}）秒。A.\(18.7\)；B.\(16.9\)；C.\(17.1\)。\par
\begin{BookAnswer}
选B。小明较慢，用时多0.9秒，小雨用时为：\[\begin{aligned}&17.8-0.9\\={}&17.8-1+0.1\\={}&16.8+0.1\\={}&16.9\text{(秒)}\end{aligned}\]
\end{BookAnswer}

% p058_s08
\BookQuestion{三、1}
直接写得数。\(20-8.5\)，\allowbreak{}\(9-0.63\)，\allowbreak{}\(7.3-2.8\)，\allowbreak{}\(1-0.36-0.64\)；\allowbreak{}\(0.77+0.23\)，\allowbreak{}\(8.88-0.8\)，\allowbreak{}\(4-1.29\)，\allowbreak{}\(6.3+0.48\)；\allowbreak{}\(2.1-0.5\)，\allowbreak{}\(4.65+5.35\)，\allowbreak{}\(6.8+2.2\)，\allowbreak{}\(18.6-8.6\)。\par
\begin{BookAnswer}
按从左到右、从上到下的顺序：\[\begin{array}{rrrr}11.5&8.37&4.5&0\\1&8.08&2.71&6.78\\1.6&10&9&10\end{array}\]其中\(1-0.36-0.64=1-(0.36+0.64)=0\)；\allowbreak{}\(0.77+0.23=1\)和\(4.65+5.35=10\)可先凑整。
\end{BookAnswer}

% p058_s09
\BookQuestion{三、2}
用竖式计算。（带★的要验算。）\par \(3.58+2.67=\textcolor{answerred}{6.25}\)\qquad \(4.6+12.95=\textcolor{answerred}{17.55}\)\par \(20.05-7.89=\textcolor{answerred}{12.16}\)\qquad ★\(50-16.34=\textcolor{answerred}{33.66}\)\par
\begin{BookAnswer}
添0的依据是小数的性质，例如\(4.6=4.60\)、\(50=50.00\)，数值没有变化。小数点对齐，空位补0；相同数位相加减。四题结果：\(6.25,17.55,12.16,33.66\)。\[\begin{array}{r}3.58\\+2.67\\\hline6.25\end{array}\quad\begin{array}{r}4.60\\+12.95\\\hline17.55\end{array}\]\[\begin{array}{r}20.05\\-7.89\\\hline12.16\end{array}\quad\begin{array}{r}50.00\\-16.34\\\hline33.66\end{array}\]最后一题用加法验算：\(33.66+16.34=50.00\)。最后一题的50先写成50.00；百分位0不足减4，要向高位逐级退位，不能把“0减4”写成4。得到的差再加减数，应还原被减数。
\end{BookAnswer}

\Needspace{8\baselineskip}

\pdfbookmark[2]{原书第59页}{orig-59}

\subsection*{原书第59页}

% p059_s01
\BookQuestion{三、3（1）}
计算下列各题，能简便计算的要简便计算。\(52.46-(8.3+12.46)\)。\par
\begin{BookAnswer}
减法性质：\(a-(b+c)=a-b-c\)；连续减去两个数时，调换这两个减数的顺序，差不变。简便法：把减去和改成连续相减，再调换两个减数的顺序。\[\begin{aligned}&52.46-(8.3+12.46)\\={}&52.46-12.46-8.3\\={}&40-8.3\\={}&31.7\end{aligned}\]
\end{BookAnswer}

% p059_s02
\BookQuestion{三、3（2）}
计算下列各题，能简便计算的要简便计算。\(86.72-10.1\)。\par
\begin{BookAnswer}
把减数10.1拆成10与0.1，先减整十数，再减0.1。\[\begin{aligned}&86.72-10.1\\={}&86.72-10-0.1\\={}&76.72-0.1\\={}&76.62\end{aligned}\]
\end{BookAnswer}

% p059_s03
\BookQuestion{三、3（3）}
计算下列各题，能简便计算的要简便计算。\(250-25.62-74.38\)。\par
\begin{BookAnswer}
把连续减去两个数改成减去它们的和，先凑成100。\[\begin{aligned}&250-25.62-74.38\\={}&250-(25.62+74.38)\\={}&250-100\\={}&150\end{aligned}\]
\end{BookAnswer}

% p059_s04
\BookQuestion{三、3（4）}
计算下列各题，能简便计算的要简便计算。\(6.9+1.3-6.9+1.3\)。\par
\begin{BookAnswer}
保留各数前的加减符号，将同样的6.9相减。\[\begin{aligned}&6.9+1.3-6.9+1.3\\={}&(6.9-6.9)+(1.3+1.3)\\={}&0+2.6\\={}&2.6\end{aligned}\]
\end{BookAnswer}

% p059_s05
\BookQuestion{四、1}
一款家用纯电动汽车满电续航\(425.8\)千米，同级燃油汽车满箱续航\(596.5\)千米。燃油汽车单次续航比新能源汽车多多少千米？\par
\begin{BookAnswer}
方法一：较长续航减去较短续航。\[\begin{aligned}&596.5-425.8\\={}&170.7\text{(千米)}\end{aligned}\]答：燃油汽车单次续航多170.7千米。\par 方法二：把减数425.8凑成426，多减0.2再加回0.2。\[\begin{aligned}&596.5-425.8\\={}&596.5-426+0.2\\={}&170.5+0.2\\={}&170.7\text{(千米)}\end{aligned}\]答：多170.7千米。
\end{BookAnswer}

% p059_s06
\BookQuestion{四、2}
亮亮家上个月节约用水\(3.65\)吨，本月积极践行低碳生活，比上月多节约用水\(1.15\)吨。两个月一共节约用水多少吨？\par
\begin{BookAnswer}
方法一：先算本月节约量，再与上月相加。\[\begin{aligned}&3.65+(3.65+1.15)\\={}&3.65+4.80\\={}&8.45\text{(吨)}\end{aligned}\]答：两个月共节约8.45吨。\par 方法二：两个月都有相当于上月的一份，本月再多1.15吨，把两份3.65先相加。\[\begin{aligned}&3.65+3.65+1.15\\={}&7.30+1.15\\={}&8.45\text{(吨)}\end{aligned}\]答：两个月共节约8.45吨。
\end{BookAnswer}

% p059_s07
\BookQuestion{四、3}
神舟飞船短期太空驻留时长为\(92.8\)天，新一代空间站组合体单次驻留可达\(186.5\)天。新一代空间站单次驻留比神舟飞船多多少天？\par
\begin{BookAnswer}
方法一：新一代驻留时长减去神舟飞船的驻留时长。\[\begin{aligned}&186.5-92.8\\={}&93.7\text{(天)}\end{aligned}\]答：新一代比神舟飞船多93.7天。\par 方法二：先减100，再把多减的7.2加回。\[\begin{aligned}&186.5-92.8\\={}&186.5-100+7.2\\={}&86.5+7.2\\={}&93.7\text{(天)}\end{aligned}\]答：多93.7天。
\end{BookAnswer}

\Needspace{8\baselineskip}

\pdfbookmark[2]{原书第60页}{orig-60}

\subsection*{原书第60页}

% p060_s01
\BookQuestion{四、4}
父亲节，琳琳用零花钱给爸爸买了一条围巾。付款100元，找回\(26.7\)元，琳琳发现多找回\(9.3\)元，围巾多少元？\par
\begin{BookAnswer}
方法一：先求原本应找回的钱，再用付款减去正确找零。\[\begin{aligned}&100-(26.7-9.3)\\={}&100-17.4\\={}&82.6\text{(元)}\end{aligned}\]答：围巾82.6元。\par 方法二：按错误找零算出的价格少了9.3元，要再加上。\[\begin{aligned}&100-26.7+9.3\\={}&73.3+9.3\\={}&82.6\text{(元)}\end{aligned}\]答：围巾82.6元。
\end{BookAnswer}

% p060_s02
\BookQuestion{四、5}
蓬莱阁景区开展文化宣传活动，准备了两种宣传画卷。长画卷长\(3.6\)米，短画卷比长画卷短\(1.35\)米。两幅画卷一共长多少米？\par
\begin{BookAnswer}
方法一：先求短画卷，再求两幅总长。\[\begin{aligned}&3.6+(3.6-1.35)\\={}&3.6+2.25\\={}&5.85\text{(米)}\end{aligned}\]答：两幅画卷共长5.85米。\par 方法二：先按两幅都长3.6米算，再减掉多算的1.35米。\[\begin{aligned}&3.6+3.6-1.35\\={}&7.2-1.35\\={}&5.85\text{(米)}\end{aligned}\]答：共长5.85米。
\end{BookAnswer}

% p060_s03
\BookQuestion{四、6（1）}
国货惠民超市促销：纯牛奶\(49.8\)元/箱，天然洗衣液\(26.35\)元/瓶，原生木浆纸巾\(13.65\)元/提。（1）购买一箱纯牛奶和一瓶天然洗衣液一共需要多少元？\par
\par\smallskip\Needspace{4\baselineskip}\begingroup\small\setlength{\tabcolsep}{3pt}
\noindent\begin{tabularx}{\linewidth}{|*{2}{>{\centering\arraybackslash}X|}}\hline
商品 & 价格\\\hline
纯牛奶 & 49.8元/箱\\\hline
天然洗衣液 & 26.35元/瓶\\\hline
原生木浆纸巾 & 13.65元/提\\\hline
\end{tabularx}\par
\endgroup\smallskip
\begin{BookAnswer}
方法一：把一箱奶价和一瓶洗衣液价相加。\[\begin{aligned}&49.8+26.35\\={}&76.15\text{(元)}\end{aligned}\]答：一共需要76.15元。\par 方法二：把49.8凑成50，多加的0.2再减去。\[\begin{aligned}&49.8+26.35\\={}&50+26.35-0.2\\={}&76.35-0.2\\={}&76.15\text{(元)}\end{aligned}\]答：一共需要76.15元。
\end{BookAnswer}

% p060_s04
\BookQuestion{四、6（2）}
国货惠民超市促销：纯牛奶\(49.8\)元/箱，天然洗衣液\(26.35\)元/瓶，原生木浆纸巾\(13.65\)元/提。（2）妈妈带100元购买这三种商品，钱够吗？若够，还剩多少元？若不够，还差多少元？\par
\par\smallskip\Needspace{4\baselineskip}\begingroup\small\setlength{\tabcolsep}{3pt}
\noindent\begin{tabularx}{\linewidth}{|*{2}{>{\centering\arraybackslash}X|}}\hline
商品 & 价格\\\hline
纯牛奶 & 49.8元/箱\\\hline
天然洗衣液 & 26.35元/瓶\\\hline
原生木浆纸巾 & 13.65元/提\\\hline
\end{tabularx}\par
\endgroup\smallskip
\begin{BookAnswer}
方法一：先求三种商品总价。\[\begin{aligned}&49.8+26.35+13.65\\={}&49.8+(26.35+13.65)\\={}&49.8+40\\={}&89.8\text{(元)}\end{aligned}\]\(89.8<100\)，所以钱够。找零：\[\begin{aligned}&100-89.8\\={}&10.2\text{(元)}\end{aligned}\]答：钱够，还剩10.2元。\par 方法二：先买奶，再扣去另两种商品的合计价。\[\begin{aligned}&100-49.8-(26.35+13.65)\\={}&50.2-40\\={}&10.2\text{(元)}\end{aligned}\]剩余为正数，说明钱够。答：还剩10.2元。
\end{BookAnswer}

\Needspace{12\baselineskip}

\section{第七单元素养评价（B）}

\Needspace{8\baselineskip}

\pdfbookmark[2]{原书第61页}{orig-61}

\subsection*{原书第61页}

% p061_s01
\BookQuestion{一、1}
整数部分最小的计数单位与小数部分最大的计数单位的和是（\hspace{1.8em}），差是（\hspace{1.8em}）。\par
\begin{BookAnswer}
整数部分最小的计数单位是1，小数部分最大的计数单位是0.1。和：\(1+0.1=1.1\)；差：\(1-0.1=0.9\)。
\end{BookAnswer}

% p061_s02
\BookQuestion{一、2}
从\(0.45\)里面去掉（\hspace{1.8em}）个\(0.01\)，结果是\(0.05\)。\par
\begin{BookAnswer}
先求去掉的数：\(0.45-0.05=0.40\)；0.40中有40个0.01，填40。
\end{BookAnswer}

% p061_s03
\BookQuestion{一、3}
比\(7.8\)少\(2.81\)的数是（\hspace{1.8em}），比\(2.36\)多\(1.4\)的数是（\hspace{1.8em}）。\par
\begin{BookAnswer}
两空依次填\(4.99\)、\(3.76\)。第一空简算：\[\begin{aligned}&7.8-2.81\\={}&7.8-2.8-0.01\\={}&5-0.01\\={}&4.99\end{aligned}\]第二空\(2.36+1.4=3.76\)。
\end{BookAnswer}

% p061_s04
\BookQuestion{一、4}
把\(7.5\)的小数点向右移动一位后，得到的数比原来多（\hspace{1.8em}）。\par
\begin{BookAnswer}
小数点向右移一位，7.5变成75；比原数多\(75-7.5=67.5\)。注意不是问移动后等于多少。
\end{BookAnswer}

% p061_s05
\BookQuestion{一、5}
在括号里填上合适的数。\((\quad)-2.4=3.58\)；\allowbreak{}\(0.86+(\quad)=1.16\)；\allowbreak{}\(4-(\quad)=0.64\)。\par
\begin{BookAnswer}
第1空：被减数=差+减数，\allowbreak{}\(3.58+2.4=5.98\)。第2空：一个加数=和−另一个加数，\allowbreak{}\(1.16-0.86=0.3\)。第3空：减数=被减数−差，\allowbreak{}\(4-0.64=3.36\)。
\end{BookAnswer}

% p061_s06
\BookQuestion{一、6（1）}
在圆圈里填上“＞”“＜”或“＝”：\(4.3-2.5\ \bigcirc\ 4.3-0.25\)。\par
\begin{BookAnswer}
被减数同为4.3，减去2.5比减去0.25多，差反而小，填“<”；也可分别算出1.8和4.05再比较。
\end{BookAnswer}

% p061_s07
\BookQuestion{一、6（2）}
在圆圈里填上“＞”“＜”或“＝”：\(3.1+9.84\ \bigcirc\ 3.84+9.1\)。\par
\begin{BookAnswer}
\(3.1+9.84=12.94\)，\allowbreak{}\(3.84+9.1=12.94\)，填“=”。也可把两数各分成整数部分和小数部分配对。
\end{BookAnswer}

% p061_s08
\BookQuestion{一、6（3）}
在圆圈里填上“＞”“＜”或“＝”：\(7.8-(3.8+1.6)\ \bigcirc\ 7.8-(3.8-1.6)\)。\par
\begin{BookAnswer}
\(3.8+1.6>3.8-1.6\)，同样用7.8去减，左边差较小，填“<”。直接计算分别为2.4和5.6。
\end{BookAnswer}

% p061_s09
\BookQuestion{一、7}
在一次跳高测试中，小明跳了\(1.25\)米，比小芳多跳了8厘米，小芳跳了（\hspace{1.8em}）米。\par
\begin{BookAnswer}
先统一长度单位，厘米化米除以100：\[8\text{厘米}=(8\div100)\text{米}=0.08\text{米}\]小芳比小明少跳0.08米，\allowbreak{}\(1.25-0.08=1.17\)米。原空后已有“米”，只填1.17。
\end{BookAnswer}

% p061_s10
\BookQuestion{一、8}
一个两位小数加上它的计数单位后是\(1.2\)，这个小数是（\hspace{1.8em}）。\par
\begin{BookAnswer}
两位小数的计数单位是0.01。原小数加0.01得1.20，所以原数为\(1.20-0.01=1.19\)。
\end{BookAnswer}

% p061_s11
\BookQuestion{一、9}
甲数是\(10.9\)，比乙数多\(0.75\)，甲、乙两数的和是（\hspace{1.8em}）。\par
\begin{BookAnswer}
甲比乙多0.75，所以乙为\(10.9-0.75=10.15\)；两数和\(10.9+10.15=21.05\)。也可算\(10.9+10.9-0.75\)。
\end{BookAnswer}

% p061_s12
\BookQuestion{一、10（1）}
按规律填数：（1）\(0.5,1.2,(\quad),2.6,(\quad)\)。\par
\begin{BookAnswer}
数列每次加0.7，依次为\(0.5,1.2,1.9,2.6,3.3\)，填1.9、3.3。
\end{BookAnswer}

% p061_s13
\BookQuestion{一、10（2）}
按规律填数：（2）\(5.19,5.16,5.13,(\quad),(\quad),5.04\)。\par
\begin{BookAnswer}
数列每次减\(0.03\)。两空依次计算：\[\begin{aligned}5.13-0.03&=5.10\\5.10-0.03&=5.07\end{aligned}\]所以填\(5.10\)、\(5.07\)；\allowbreak{}\(5.10=5.1\)。
\end{BookAnswer}

% p061_s14
\BookQuestion{一、11（两空，续下页）}
天天家、云云家、学校在同一条直线上。天天家到学校的距离是\(1.36\)千米，云云家到学校的距离是\(0.78\)千米。天天家和云云家最远相距（\hspace{1.8em}）千米，最近相距（\hspace{1.8em}）千米。（后一空续下页。）\par
\begin{BookAnswer}
学校在两家之间时，两家相距最远，\allowbreak{}\(1.36+0.78=2.14\)千米；两家在学校同一侧时，相距最近，\allowbreak{}\(1.36-0.78=0.58\)千米。不能只画一种位置关系。可先仿照线段示意图分别画“学校在两家之间”和“两家在学校同一侧”，再选择求和或求差，避免漏掉一种位置。
\end{BookAnswer}

\Needspace{8\baselineskip}

\pdfbookmark[2]{原书第62页}{orig-62}

\subsection*{原书第62页}

% p062_s01
\BookQuestion{一、11（续）}
承原书61页第11题：天天家、云云家、学校在同一条直线上，天天家距学校\(1.36\)千米，云云家距学校\(0.78\)千米。两家最近相距（\hspace{1.8em}）千米。\par
\begin{BookAnswer}
最近距离填0.58千米。两家在学校同侧，两段到校路程重合，求差：\(1.36-0.78=0.58\)。最远2.14千米已填在上一页。
\BookFigure{figures/p062_s01-01.tex}
\end{BookAnswer}

% p062_s02
\BookQuestion{一、12}
一桶油连桶重\(46.8\)千克，倒出一半油后连桶重\(24.5\)千克，桶重（\hspace{1.8em}）千克，原来有（\hspace{1.8em}）千克油。\par
\begin{BookAnswer}
前后相差的质量就是半桶油的质量：\[\begin{aligned}&46.8-24.5\\={}&22.3\text{(千克)}\end{aligned}\]原油质量：\[\begin{aligned}&22.3+22.3\\={}&44.6\text{(千克)}\end{aligned}\]桶的质量：\[\begin{aligned}&46.8-44.6\\={}&2.2\text{(千克)}\end{aligned}\]所以两空依次填2.2、44.6。另一思路：两份“半桶油连桶”的质量，比一份“整桶油连桶”的质量多一个桶。\[\begin{aligned}&24.5+24.5-46.8\\={}&49-46.8\\={}&2.2\text{(千克)}\end{aligned}\]再算\(46.8-2.2=44.6\)(千克)。
\end{BookAnswer}

% p062_s03
\BookQuestion{二、1}
算式\(2\square.4+9.\square6\)的结果可能是（\hspace{1.8em}）。A.\(29.50\)；B.\(29.36\)；C.\(31.26\)。\par
\begin{BookAnswer}
两个加数的小数部分相加，百分位必为6，不可能得29.50；最小和\(20.4+9.06=29.46\)，也不可能得29.36。\(21.4+9.86=31.26\)确实可能，选C。
\end{BookAnswer}

% p062_s04
\BookQuestion{二、2}
在右图的加法竖式中，虚线框中的“9”和“5”相加得（\hspace{1.8em}）。A.14个一；B.14个十分之一；C.14个百分之一。竖式为\(34.95+2.5\)，小数点对齐。\par
\BookFigure{figures/p062_s04-01.tex}
\begin{BookAnswer}
9与5在十分位，表示9个十分之一和5个十分之一，合为14个十分之一，即1.4。选B。
\end{BookAnswer}

% p062_s05
\BookQuestion{二、3}
与\(15.8+9.9\)计算结果不同的是（\hspace{1.8em}）。A.\(15.8+9+0.9\)；B.\(15.8+10+0.1\)；C.\(15.8+10-0.1\)。\par
\begin{BookAnswer}
\(9.9=9+0.9=10-0.1\)。A与C都等于\(15.8+9.9\)；B为\(15.8+10+0.1\)，比正确值多0.2，选B。
\end{BookAnswer}

% p062_s06
\BookQuestion{二、4}
两个数相加，如果一个数增加\(3.57\)，另一个数减少\(4.5\)，和（\hspace{1.8em}）。A.减少\(1.07\)；B.增加\(0.93\)；C.减少\(0.93\)。\par
\begin{BookAnswer}
一个加数增加3.57，另一个减少4.5，少掉的比多出的多\(4.5-3.57=0.93\)，即和减少0.93，选C。
\end{BookAnswer}

% p062_s07
\BookQuestion{二、5}
下面的数最接近9的是（\hspace{1.8em}）。A.\(9.02\)；B.\(8.99\)；C.\(9.1\)。\par
\begin{BookAnswer}
与9分别相差\(9.02-9=0.02\)、\(9-8.99=0.01\)、\(9.1-9=0.1\)。差最小的是8.99，选B。
\end{BookAnswer}

% p062_s08
\BookQuestion{二、6}
下列说法中，正确的有（\hspace{1.8em}）个。①两个数的和是小数，这两个数也一定都是小数。②\(6.4+3.6-6.4+3.6=0\)。③计算小数加减法，应把末位对齐。④在减法算式中，被减数和减数同时增加\(2.5\)，差不变。A.0；B.1；C.2。\par
\begin{BookAnswer}
①错：如\(1+0.2=1.2\)，和是小数但加数不全是小数。②错：\(6.4+3.6-6.4+3.6=7.2\)。③错：应对齐小数点，不是只对齐末位。④对：被减数、减数同时增加相同的数，差不变。仅1个正确，选B。
\end{BookAnswer}

% p062_s09
\BookQuestion{三、1}
直接写得数。\(3.25+5.5\)，\allowbreak{}\(1.7+2.5\)，\allowbreak{}\(6.76-0.6\)，\allowbreak{}\(2.6+8\)；\allowbreak{}\(3.8+0.12\)，\allowbreak{}\(2.36+1.64\)，\allowbreak{}\(4.2-4\)，\allowbreak{}\(10-4.5\)。\par
\begin{BookAnswer}
按行依次为\[\begin{array}{rrrr}8.75&4.2&6.16&10.6\\3.92&4&0.2&5.5\end{array}\]小数点对齐；\allowbreak{}\(2.36+1.64\)可先凑成4；\allowbreak{}\(4.2-4=0.2\)，整数部分是0。
\end{BookAnswer}

\Needspace{8\baselineskip}

\pdfbookmark[2]{原书第63页}{orig-63}

\subsection*{原书第63页}

% p063_s01
\BookQuestion{三、2}
用竖式计算。（带★的要验算。）\par ★\(8.7+12.36=\textcolor{answerred}{21.06}\)\qquad \(5.1-3.59=\textcolor{answerred}{1.51}\)\par ★\(12-2.25=\textcolor{answerred}{9.75}\)\par
\begin{BookAnswer}
相同数位对齐，必要时补0。\[\begin{array}{r}8.70\\+12.36\\\hline21.06\end{array}\quad\begin{array}{r}5.10\\-3.59\\\hline1.51\end{array}\quad\begin{array}{r}12.00\\-2.25\\\hline9.75\end{array}\]带星号两题要验算：\(21.06-12.36=8.70\)；\allowbreak{}\(9.75+2.25=12.00\)。加法可用减法验算，减法可用加法验算。\(5.1=5.10\)、\(12=12.00\)来自小数末尾添0的性质。退位时借来的1个十分之一相当于10个百分之一；12.00连续退位要逐级处理。
\end{BookAnswer}

% p063_s02
\BookQuestion{三、3（1）}
计算下列各题，能简便计算的要简便计算。\(35-3.7+6.3\)。\par
\begin{BookAnswer}
先算6.3与3.7的差，再与整数35相加。\[\begin{aligned}&35-3.7+6.3\\={}&35+(6.3-3.7)\\={}&35+2.6\\={}&37.6\end{aligned}\]
\end{BookAnswer}

% p063_s03
\BookQuestion{三、3（2）}
计算下列各题，能简便计算的要简便计算。\(27-9.98\)。\par
\begin{BookAnswer}
把9.98凑成10，多减的0.02再加回。\[\begin{aligned}&27-9.98\\={}&27-10+0.02\\={}&17+0.02\\={}&17.02\end{aligned}\]
\end{BookAnswer}

% p063_s04
\BookQuestion{三、3（3）}
计算下列各题，能简便计算的要简便计算。\(45.35-(7.2+5.35)\)。\par
\begin{BookAnswer}
先减去5.35，可以得到整数40。\[\begin{aligned}&45.35-(7.2+5.35)\\={}&45.35-5.35-7.2\\={}&40-7.2\\={}&32.8\end{aligned}\]
\end{BookAnswer}

% p063_s05
\BookQuestion{三、3（4）}
计算下列各题，能简便计算的要简便计算。\(23.8-3.9-5.1+36.2\)。\par
\begin{BookAnswer}
运算律和减法性质同样适用于小数，但移动一个数时应连同它前面的加减符号一起移动。两个加数先凑60，两个减数先凑9。\[\begin{aligned}&23.8-3.9-5.1+36.2\\={}&(23.8+36.2)-(3.9+5.1)\\={}&60-9\\={}&51\end{aligned}\]
\end{BookAnswer}

% p063_s06
\BookQuestion{四、1}
2024年4月30日，神舟十七号成功返航，用时约9小时，比神舟十二号的返航时间缩短了约\(19.5\)小时。神舟十二号返航大约用了多少小时？\par
\begin{BookAnswer}
方法一：神舟十七号比十二号用时短19.5小时，求十二号用加法。\[\begin{aligned}&9+19.5\\={}&28.5\text{(小时)}\end{aligned}\]答：神舟十二号返航大约用了28.5小时。\par 方法二：把19.5凑成20，再减去多加的0.5。\[\begin{aligned}&9+19.5\\={}&9+20-0.5\\={}&29-0.5\\={}&28.5\text{(小时)}\end{aligned}\]答：大约用了28.5小时。
\end{BookAnswer}

% p063_s07
\BookQuestion{四、2}
把一根长2米的竹竿竖直插入水池中，竹竿入泥部分长16厘米，露出水面部分长54厘米。水池的水深多少米？\par
\begin{BookAnswer}
竹竿总长减入泥和露出两段，才是水深。先统一成米：\[\begin{aligned}16\text{厘米}&=0.16\text{米}\\54\text{厘米}&=0.54\text{米}\end{aligned}\]利用连减性质先凑0.7：\[\begin{aligned}&2-0.16-0.54\\={}&2-(0.16+0.54)\\={}&2-0.7\\={}&1.3\text{(米)}\end{aligned}\]也可统一成厘米：\[2\text{米}=200\text{厘米}\]\[\begin{aligned}&200-(16+54)\\={}&200-70\\={}&130\text{(厘米)}\end{aligned}\]\[130\text{厘米}=1.3\text{米}\]答：水池的水深是1.3米。
\end{BookAnswer}

\Needspace{8\baselineskip}

\pdfbookmark[2]{原书第64页}{orig-64}

\subsection*{原书第64页}

% p064_s01
\BookQuestion{四、3}
“五岳”是我国五大名山的总称，其中海拔最低的是南岳衡山，海拔约\(1.3\)千米，海拔最高的是西岳华山，西岳华山的海拔比南岳衡山约高\(0.85\)千米，比东岳泰山约高\(0.62\)千米，东岳泰山约高多少千米？\par
\begin{BookAnswer}
方法一：先求华山海拔，再减去华山比泰山高的部分。\[\begin{aligned}&1.3+0.85-0.62\\={}&2.15-0.62\\={}&1.53\text{(千米)}\end{aligned}\]答：东岳泰山约高1.53千米。\par 方法二：先求泰山比衡山高多少，再加衡山海拔。\[\begin{aligned}&1.3+(0.85-0.62)\\={}&1.3+0.23\\={}&1.53\text{(千米)}\end{aligned}\]答：约高1.53千米。
\end{BookAnswer}

% p064_s02
\BookQuestion{四、4}
一根绳子长\(2.35\)米，另一根绳子长\(2.85\)米，把两根绳子接在一起后的长度是\(4.7\)米，接头处是多少米？\par
\begin{BookAnswer}
方法一：两根原长的和减去接后总长，就是重复计算的接头长。\[\begin{aligned}&2.35+2.85-4.7\\={}&5.20-4.70\\={}&0.50\text{(米)}\end{aligned}\]答：接头处长0.5米。\par 方法二：先从接后总长中扣去第二根，求第一根未重叠部分，再与其原长比较。\[\begin{aligned}&2.35-(4.7-2.85)\\={}&2.35-1.85\\={}&0.50\text{(米)}\end{aligned}\]答：接头处长0.5米。
\end{BookAnswer}

% p064_s03
\BookQuestion{四、5}
电商平台助力乡村振兴，李阿姨通过直播销售家乡农产品。9月份售出\(1.8\)万箱农产品，10月份比9月份少\(0.26\)万箱，这两个月一共售出多少万箱？\par
\begin{BookAnswer}
方法一：先求10月销量，再加9月销量。\[\begin{aligned}&1.8+(1.8-0.26)\\={}&1.8+1.54\\={}&3.34\text{(万箱)}\end{aligned}\]答：两个月共售出3.34万箱。\par 方法二：按两个月都卖1.8万箱算，再减去10月少卖的0.26万箱。\[\begin{aligned}&1.8+1.8-0.26\\={}&3.60-0.26\\={}&3.34\text{(万箱)}\end{aligned}\]答：共售出3.34万箱。
\end{BookAnswer}

% p064_s04
\BookQuestion{四、6}
小明去书店买了一本书，付给售货员50元，找回\(38.5\)元。他仔细一算，发现售货员多找了10元，这本书多少元？\par
\begin{BookAnswer}
方法一：先求正确找零，再用付款减去正确找零。\[\begin{aligned}&50-(38.5-10)\\={}&50-28.5\\={}&21.5\text{(元)}\end{aligned}\]答：这本书21.5元。\par 方法二：按错误找零算出的书价少了10元，再加回10元。\[\begin{aligned}&50-38.5+10\\={}&11.5+10\\={}&21.5\text{(元)}\end{aligned}\]答：这本书21.5元。
\end{BookAnswer}
